
\documentclass[12pt,a4paper]{article}

% 引入必要的宏包
\usepackage{ctex} % 支持中文排版
\usepackage{amsmath, amssymb} % 数学公式和符号支持
\usepackage{geometry} % 页面边距设置
\geometry{left=2.5cm, right=2.5cm, top=2.5cm, bottom=2.5cm}
\usepackage{hyperref} % 支持书签和链接
\hypersetup{colorlinks=true, linkcolor=blue}

\usepackage{xcolor}
\usepackage{eso-pic}
\usepackage{graphicx}

\newcommand{\WMText}{贺昱超学习资料}
\newcommand{\WMScale}{2.28}
\newcommand{\WMAngle}{45}
\colorlet{WMColor}{black!8}

\AddToShipoutPictureBG{%
  \AtPageLowerLeft{%
    \put(\LenToUnit{0.66\paperwidth},\LenToUnit{0.70\paperheight}){%
      \rotatebox{\WMAngle}{\scalebox{\WMScale}{\textcolor{WMColor}{\WMText}}}%
    }%
    \put(\LenToUnit{0.33\paperwidth},\LenToUnit{0.50\paperheight}){%
      \rotatebox{\WMAngle}{\scalebox{\WMScale}{\textcolor{WMColor}{\WMText}}}%
    }%
    \put(\LenToUnit{0.66\paperwidth},\LenToUnit{0.30\paperheight}){%
      \rotatebox{\WMAngle}{\scalebox{\WMScale}{\textcolor{WMColor}{\WMText}}}%
    }%
  }%
}

\title{\textbf{微积分与 CCL 竞赛大纲}}
\author{}
\date{}

\begin{document}

\maketitle

\section*{模块一：微积分知识体系 }

\subsection*{1. 极限与连续 (Limits and Continuity)}
\begin{itemize}
    \item \textbf{极限的概念：} 左极限、右极限、双侧极限、无穷极限、极限在无穷远处。
    \item \textbf{极限的计算：} 代数化简、夹逼定理 (Squeeze Theorem)、洛必达法则 (L'H\^opital's Rule)。
    \item \textbf{连续性：} 连续的三要素、各类间断点（可去、跳跃、无穷）。
    \item \textbf{介值定理 (Intermediate Value Theorem, IVT)：} 连续函数在区间内的取值特性。
\end{itemize}

\subsection*{2. 导数基础 (Derivatives and Differentiation)}
\begin{itemize}
    \item \textbf{定义：} 瞬时变化率、割线到切线的极限过程（导数的极限定义）。
    \item \textbf{运算规则：} 幂法则、指数/对数导数、三角/反三角函数导数、乘积法则、商法则、链式法则。
    \item \textbf{高级技巧：} 隐函数求导、对数求导法、反函数求导法则。
\end{itemize}

\subsection*{3. 导数的应用 (Applications of Differentiation)}
\begin{itemize}
    \item \textbf{几何意义：} 切线与法线方程、局部线性近似 (Linearization)。
    \item \textbf{核心定理：} 极值定理 (EVT)、罗尔定理 (Rolle's Theorem)、均值定理 (MVT)。
    \item \textbf{函数图象分析：} 临界点、极大/极小值（一阶求导测试）、拐点与凹凸性（二阶求导测试）。
    \item \textbf{实际应用：} 相关变化率 (Related Rates)、绝对极值与最优化问题 (Optimization)、物理运动学（位置、速度、加速度之间的关系）。
\end{itemize}

\subsection*{4. 积分基础 (Integration)}
\begin{itemize}
    \item \textbf{不定积分：} 反导数概念、基本积分法则表。
    \item \textbf{定积分：} 黎曼和 (Riemann Sums，如左/右/中点/梯形法则)、面积的极限定义。
    \item \textbf{微积分基本定理 (FTC)：}
    \begin{itemize}
        \item 第一部分：联系导数与积分（求解变上限积分函数的导数）。
        \item 第二部分：利用反导数计算定积分（牛顿-莱布尼茨公式）。
    \end{itemize}
    \item \textbf{积分技巧：} 换元积分法 (u-Substitution)、分部积分法 (Integration by Parts)、部分分式展开法 (Partial Fractions)、三角换元。
    \item \textbf{反常积分 (Improper Integrals)：} 包含无穷边界或无界函数的积分收敛性问题。
\end{itemize}

\subsection*{5. 积分的应用 (Applications of Integration)}
\begin{itemize}
    \item \textbf{几何应用：} 曲线下方的面积、两条或者多条曲线围成的面积。
    \item \textbf{体积计算：} 旋转体体积（圆盘法 Disk Method、垫圈法 Washer Method、圆柱壳法 Shell Method）、已知截面面积的立体体积。
    \item \textbf{其他计算：} 曲线弧长 (Arc Length)、旋转体表面积。
    \item \textbf{物理应用：} 函数的平均值定理、变力做功 (Work)、流体静压力、质心计算 (Center of Mass)。
\end{itemize}

\subsection*{6. 微分方程与级数 (Differential Equations \& Infinite Series)}
\begin{itemize}
    \item \textbf{微分方程：} 斜率场 (Slope Fields)、欧拉方法 (Euler's Method)、求解可分离变量的微分方程、指数与逻辑斯蒂增长模型 (Logistic Growth)。
    \item \textbf{参数与极坐标：} 参数方程求导与积分、极坐标系统下的微积分计算。
    \item \textbf{无穷级数：} 级数收敛性测试、泰勒级数与麦克劳林级数 (Taylor and Maclaurin Series)、幂级数及其收敛区间与收敛半径、误差估算。
\end{itemize}

\vspace{1em}
\hrule
\vspace{1em}

\section*{模块二：CCL大纲}
结合 2025 CCL Calculus 试卷的题型，以下是该竞赛重点考察的具体知识点和应用方向。试卷整体侧重于计算的熟练度、概念的灵活转化以及微积分在物理和几何中的综合应用。

\subsection*{1. 极限与连续的进阶计算}
\begin{itemize}
    \item \textbf{带绝对值的单侧极限：} 需要判断去绝对值符号时的正负号，例如计算 $\lim_{x \to 3^{-}} \frac{9-x^2}{|x-3|}$。
    \item \textbf{分段函数极限：} 要求熟练求解分段函数在分界点处的左极限与右极限。
    \item \textbf{导数的极限逆向识别：} 给定复杂的极限表达式，能识别出它是某个特定函数在某一点的导数。例如识别对称差商 $\lim_{h \to 0} \frac{f(2+h)-f(2-h)}{h} = 2f'(2)$。
    \item \textbf{水平渐近线：} 快速比较分子分母多项式的最高次幂，求解 $x \to \infty$ 时的极限。
\end{itemize}

\subsection*{2. 导数与微分的综合运算}
\begin{itemize}
    \item \textbf{复合与嵌套法则：} 熟练结合商法则、乘积法则与链式法则。例如对 $y = \frac{\sin x}{e^x}$ 求导。
    \item \textbf{隐函数求导：} 对非显式方程直接进行等式两边求导。例如对方程 $e^y = x+y^2$ 求导，并进一步求解二阶导数 $\frac{d^2y}{dx^2}$。
    \item \textbf{非常规变量求导：} 理解导数的本质是变化率之比。例如计算 $\frac{dy}{d(\sqrt{x+4})}$ 而不是常规的 $\frac{dy}{dx}$。
\end{itemize}

\subsection*{3. 导数的几何与物理应用}
\begin{itemize}
    \item \textbf{线性近似 (Linear Approximation)：} 利用已知点的函数值和导数值，估算附近点的函数值。例如利用 $f(5)$ 和 $f'(5)$ 估算 $f(5.02)$。
    \item \textbf{图象特征判定：} 将“递增/递减”与一阶导数的正负挂钩，将“变化率的增加/减少”与二阶导数的正负（凹凸性）挂钩。
    \item \textbf{均值定理 (MVT) 的代数求解：} 在闭区间上应用中值定理，建立方程并求解满足条件的未知参数 $c$ 或 $a$。
    \item \textbf{相关变化率 (Related Rates)：} 建立多变量之间的代数关系，两边对时间 $t$ 求导。例如已知 $x$ 和 $y$ 的变化率，求 $z = \frac{x^2y}{x+y}$ 的变化率。
    \item \textbf{最优化设计：} 结合几何图形寻找极值（如面积固定时求最小周长）。
    \item \textbf{运动学分析：} 在直线运动和抛体运动中，熟练在位置 $s(t)$、速度 $v(t)$ 和加速度 $a(t)$ 之间进行求导或积分转换。例如已知加速度 $a = -32$，结合初始条件求抛体运动的最大高度。
\end{itemize}

\subsection*{4. 积分计算与微分方程}
\begin{itemize}
    \item \textbf{微积分基本定理的应用：} 针对变上限积分函数能够直接求导并代入复杂边界。例如已知 $F(x) = \int_{1}^{x} \frac{3}{4+\ln t} dt$，求 $F'(e^2)$。
    \item \textbf{绝对值函数的积分：} 在计算两条曲线围成的面积时，准确处理含有绝对值的函数（如 $y = |x+8|$），并找到正确的积分区间与交点。
    \item \textbf{积分均值定理：} 利用公式 $\frac{1}{b-a} \int_{a}^{b} f(x) dx$ 求解函数在某区间上的平均值。
    \item \textbf{可分离变量微分方程：} 根据实际物理情境列出微分方程并求解。例如水箱漏水问题中的 $\frac{dw}{dt} = -k\sqrt{w}$，分离变量并积分求解特定时间点。
    \item \textbf{曲线族的推导：} 根据给定的切线斜率特征列出微分方程并求解。例如斜率满足 $y' = -\frac{x}{y}$ 时，积分求解出该曲线族代表的几何图形。
\end{itemize}

\end{document}

```